% =====================================================================
%  22  付録 コスト推算 3 — OPEX 運転費の推算と単価・TAC
%  source_memo: 03_economic_indicators.tex §3.3 ＋ flowsheet/economics.py
%   （apply_hasebe_aggregation / _compute_labor_cost / collect_capex_opex）
%   ＋ src/cost_parameters.py（単価）＋ utility_selector.py（用役階層）。
%   長谷部・外輪「プロセス設計 No.3」§3.3/§3.4 準拠。
%  方針：OPEX 経験式と単価をシミュレータの実装どおり提示。丸括弧禁止・最小色。
% =====================================================================
\begin{frame}[t]

  \vspace*{1.2mm}
  {\footnotesize 年間運転費 OPEX は\textbf{長谷部・外輪の経験式}で推算。原料・用役の実費に、
   保全・労務・諸経費を経験係数で上乗せする。}\par

  \vspace{0.5em}
  {\small\boldmath$\mathrm{OPEX}=0.180\,C_{TM}+2.73\,C_{OL}+1.23\,(C_{UT}+C_{WT}+C_{RM})$\unboldmath}\par
  \vspace{0.2em}
  {\scriptsize $C_{TM}$ 建設費 ｜ $C_{OL}$ 労務費 ｜ $C_{UT}$ 用役費 ｜ $C_{WT}$ 廃棄物費 $=0$ ｜ $C_{RM}$ 原料費\\
   {\color{black!60}＋ 触媒・吸着剤・膜の交換費を別建て}}\par

  \vspace{0.7em}
  \begin{columns}[T,onlytextwidth]
    % ===== 左：労務費の式 ＋ 別建て交換費 =====
    \begin{column}{0.50\textwidth}
      {\scriptsize\bfseries\color{HeaderBlue} 労務費 $C_{OL}$　主要機器数から運転員数を決定}\par
      \vspace{0.2em}
      {\scriptsize $N_{OL}=\big\lceil\sqrt{6.29+0.21\,N_{\mathrm{eq}}}\,\big\rceil$\quad $1$ 班の運転員数}\par
      \vspace{0.15em}
      {\scriptsize $4$ 直 $\to$ 総員 $=4\,N_{OL}$、年俸 $600$ 万円/人\par
       $N_{\mathrm{eq}}$：反応器・塔・圧縮機・熱交換器などの数}\par

      \vspace{0.7em}
      {\scriptsize\bfseries\color{HeaderBlue} 別建て交換費　経験式の枠外・$1.23$ 倍なし}\par
      \vspace{0.2em}
      {\scriptsize\setlength{\tabcolsep}{4pt}\renewcommand{\arraystretch}{1.2}%
       \begin{tabular}{@{}l l@{}}
         触媒 $\mathrm{Cr_2O_3}$ & $23$ USD/kg $\div$ 寿命 $2.5$ 年 \\
         PSA 活性炭 & $5$ USD/kg $\div$ $4$ 年 \\
         膜モジュール & 膜本体 CAPEX $\div$ $3$ 年 \\
       \end{tabular}}\par
    \end{column}

    \begin{column}{0.02\textwidth}\end{column}

    % ===== 右：単価一覧 =====
    \begin{column}{0.47\textwidth}
      {\scriptsize\bfseries\color{HeaderBlue} 原料・製品・用役の単価}\par
      \vspace{0.2em}
      {\setlength{\fboxsep}{4pt}\setlength{\fboxrule}{0.8pt}%
       \fcolorbox{HeaderBlue}{white}{%
        \begin{minipage}{\dimexpr\linewidth-2\fboxsep-2\fboxrule\relax}
        {\scriptsize\setlength{\tabcolsep}{3pt}\renewcommand{\arraystretch}{1.18}%
         \begin{tabular}{@{}l r@{}}
           原料 LPG & $95$ 円/kg \\
           製品 $\mathrm{C_3H_6}$ & {\color{SpRed}$150$} 円/kg \\
           製品 $\mathrm{H_2}$ & {\color{SpRed}$400$} 円/kg \\
           \midrule
           電力 & $17$ 円/kWh \\
           蒸気 LP/MP/HP & $1050$/$1120$/$1330$ 円/GJ \\
           燃料 & $1830$ 円/GJ \\
           冷媒 & Carnot 動力で算定 \\
         \end{tabular}}
        \end{minipage}}}\par
      \vspace{0.18em}
      {\tiny\color{black!55} 用役単価は海外 TEA 文献に CEPCI・為替補正。稼働 $8000$ h/年}\par
    \end{column}
  \end{columns}

  \vspace{0.8em}
  % ===== 最下部：TAC のまとめ帯 =====
  {\setlength{\fboxsep}{5pt}%
   \noindent\colorbox{HiBlue}{%
    \begin{minipage}{\dimexpr\linewidth-2\fboxsep\relax}
     \centering\small\boldmath
     {\color{SpRed}\textbf{TAC}}$\ =\ \mathrm{CAPEX}\,/\,8\ \text{年}\ +\ \mathrm{OPEX}$
     \unboldmath{\footnotesize\quad 建設費を $8$ 年で回収した年資本費＋年運転費}
    \end{minipage}}}

\end{frame}
